% Source: 02 - Tue Sep 29 2026.pdf, page 11. Content remains in source order.
\sourcepage{02}{11}

\subsection{Formal languages and concrete decision problems}
There is a one-to-one correspondence between concrete problems and formal languages over $\Sigma=\{0,1\}$:
\[
\Pi_C\ \longleftrightarrow\ L_{\Pi_C}=\{x\in\bits^*:\Pi_C(x)=1\}
\quad\text{(encodings of positive instances)}.
\]
We will use the two concepts interchangeably.
\[
\begin{aligned}
\text{Problem }\Pi_C&\longleftrightarrow\text{Algorithm solving }\Pi_C,\\
\text{Language }L_{\Pi_C}&\longleftrightarrow\text{Algorithm deciding }L_{\Pi_C}\ \text{(Automaton)}.
\end{aligned}
\]

\textbf{DEF:}
\[
L_A=\{x\in\bits^*:\ A(x)=1\}\quad\text{(language decided by }A\text{)}.
\]

\begin{itemize}
    \item An automaton $A$ decides a language $L$ if it accepts all strings in $L$ and rejects all strings not in $L$.
    \item An algorithm that solves a concrete decision problem $\Pi_C$ is an algorithm that decides the language $L_{\Pi_C}$. It returns 1 for all strings in $L_{\Pi_C}$ and 0 for all strings not in $L_{\Pi_C}$.
    \item Algorithm $A$ accepts $x$ if $A(x)=1$. $A$ rejects $x$ if $A(x)=0$.
\end{itemize}

